\chapter{数据质量与证据可信度}
\label{ch:week04-data-quality}

指标计算依赖记录的覆盖范围、测量方式和时间含义。订单漏录会改变统计对象，日期
补录会改变历史可用信息，单位变更会改变数值的解释。数据质量分析研究这些记录
特征是否满足既定用途的要求，以及记录误差如何影响统计量和决策。

质量要求由数据用途确定，数据剖析描述记录的结构与分布，缺失、测量误差和比率
误差则影响统计结论。业务要求可转化为检查规则，修复与发布策略需依据预计损失
选择（第~\ref{ch:week01-intro}~章）。质量分析同时规定统计推断的输入条件。

\section{输入数据与分析结果}

第~\ref{ch:week03-sql-kpi}~章的零售样例中，已交付订单准时率可以写为
\[
\mathrm{OnTimeRate}_{\rm delivered}
=
\frac{\sum_i I(a_i\le p_i)}{n},
\]
其中 \(a_i\) 是实际交付日期，\(p_i\) 是承诺交付日期，\(n\) 是被纳入统计的已交付订单数。该计算依赖以下数据质量假设：
\begin{itemize}
  \item 每个订单有唯一且稳定的 \texttt{order\_id}；
  \item 每个订单能连接到存在的客户和产品；
  \item \texttt{actual\_date} 和 \texttt{promised\_date} 是可比较的日期；
  \item 已交付样本与已到期承诺样本已区分，未完成订单按所选口径处理；
  \item 时间窗口使用订单日期、承诺日期还是实际日期已有明确约定；
  \item 指标刷新时使用的数据批次、计算代码和业务口径可以被追溯。
\end{itemize}

重复订单会放大销售额，缺失外键可能造成连接丢行，日期错误会改变准时判断。检查规则需要按指标用途确定。已交付订单准时率以已完成交付为纳入条件，到期承诺口径则保留到期未完成的订单。

\begin{definition}[指标可信度]
一个管理指标具有可信度，是指其事实来源、对象范围、质量假设、计算规则、失败记录、刷新时间和责任边界明确，并且可根据来源记录复核。
\end{definition}

指标的计算依据按以下过程形成。
\[
\text{业务过程}
\rightarrow
\text{数据生成}
\rightarrow
\text{质量规则}
\rightarrow
\text{指标计算}
\rightarrow
\text{管理解释}
\rightarrow
\text{行动复盘}.
\]
质量规则规定数据生成与指标计算之间的可检验条件。

\section{数据生成、质量维度与统计误差}
\subsection{数据质量的基本定义}
\subjectindex{shujuzhiliang}{数据质量}

数据质量依据用途定义。同一份数据可能满足一种用途的质量要求，同时无法满足另一种用途的要求。例如，月度销售概览可能可以容忍少量非核心字段缺失，但客户 SLA 违约判断不能容忍战略客户订单的交付日期缺失；趋势分析可以接受一两天的数据入仓延迟，但实时补货调度不能接受这样的延迟。

管理决策中常用的工作定义是：数据是否适合当前业务、分析和行动用途。

ISO/IEC 25012 将数据质量模型用于结构化数据的质量要求、度量和评估；Wang and Strong 从数据消费者角度提出内在质量、情境质量、表征质量和可访问性等信息质量类别；DAMA-DMBOK 则把数据质量放在数据治理、元数据、主数据、数据仓库和商业智能等知识域之间。

评价数据质量需要先说明用途，再据此选择质量维度与度量方法。规则符合情况提供
可计算的检查结果，来源和处理记录提供核查依据。例如，研究交付周期时，要考察
订单覆盖是否完整、日期是否准确、时间单位是否一致，以及补录记录是否可识别。

\begin{definition}[用途适配的数据质量]
给定业务问题 \(B\)、指标或模型 \(M\)、数据集 \(D\) 和质量规则集合 \(\Phi\)，若 \(D\) 在对象范围、字段含义、时间窗口、取值约束、连接关系与来源记录等方面足以支撑 \(M\) 回答 \(B\)，则称 \(D\) 对该用途具有足够的数据质量。
\end{definition}

同一数据集用于不同任务时，可以有不同的质量要求。月度收入分析关注交易覆盖与金额准确性，实时调度还要求状态更新及时。质量评价因而须与具体用途和容许误差联系。

\subsection{数据生成过程与质量风险}
\subjectindex{shujushengchengguocheng}{数据生成过程}

企业订单数据从业务发生到指标发布经历多个环节，各环节均可能引入误差。
\[
\begin{aligned}
\text{客户下单}
&\rightarrow
\text{业务系统录入}
\rightarrow
\text{主数据匹配}
\rightarrow
\text{履约状态更新}\\
&\rightarrow
\text{数据抽取与入仓}
\rightarrow
\text{指标计算}.
\end{aligned}
\]
人工录入会产生空值和拼写错误；系统迁移会产生编码不一致；接口延迟会造成数据不及时；主数据缺失会造成外键无法匹配；指标口径变化若未版本化，会导致同名指标在不同月份不可比较。

\begin{table}[htbp]
\centering
\caption{数据生成环节与质量风险}
\begin{tabular}{P{0.22\linewidth}P{0.35\linewidth}P{0.30\linewidth}}
\toprule
\rowcolor{BDMTableHead}
环节 & 常见风险 & 管理后果\\
\midrule
业务录入 & 必填字段为空、数量录入错误。 & 订单无法归属或金额失真。\\
主数据维护 & 客户、产品、渠道编码不稳定。 & 分组比较口径变化。\\
系统集成 & 接口延迟、字段截断、重复发送。 & 数据不及时或重复计数。\\
仓库建模 & 事实粒度混乱、连接键不唯一。 & KPI 分子或分母错误。\\
指标发布 & 口径变更未版本化。 & 管理会议中出现无法解释的差异。\\
\bottomrule
\end{tabular}
\end{table}

\begin{remark}
质量事件响应需定位产生误差的环节，确定受影响的指标与修复责任，并决定发布控制、业务流程及数据契约的调整方式。
\end{remark}

\subsection{规则符合性与记录准确性}
规则检查判断记录是否满足给定条件，准确性检查比较记录与实际业务。假设某订单
实际在下单后第二十天交付，却被录为第十天，该记录仍满足日期格式、先后顺序与
非空规则。确认日期准确性需要与签收凭证或其他独立来源核对。

异常记录的处理取决于业务含义。负数量可以表示退货，缺少发货日期可以表示订单
尚未发运。核对状态以后，再决定更正、补录、保留例外或修改规则，并记录原值、
处理依据和处置结果。

质量体系通常结合三类证据：结构规则检查内部一致性，对账将系统记录与独立来源
比较，抽样复核检验记录是否符合实际业务。各类证据的覆盖范围均有限。
结构规则便于规模化运行，对账依赖参照来源本身可信，人工复核则受成本和样本选择
限制。特定行动所需的证据组合应结合这些限制确定。

\subsection{数据质量维度}

常见数据质量维度包括完整性、唯一性、有效性、一致性、及时性、准确性和可追溯性。

\begin{table}[htbp]
\centering
\caption{常见数据质量维度}
\begin{tabular}{P{0.15\linewidth}P{0.28\linewidth}P{0.28\linewidth}P{0.18\linewidth}}
\toprule
\rowcolor{BDMTableHead}
维度 & 基本问题 & 订单数据示例 & 典型测试\\
\midrule

完整性 & 必要字段是否存在？ & 订单缺少客户编号或实际交付日期。 & not null、缺失率。\\
唯一性 & 业务对象是否被重复记录？ & 同一 \texttt{order\_id} 出现多行。 & unique、重复键。\\
有效性 & 字段是否满足取值域和格式？ & 渠道不属于 Direct、Partner、Online。 & accepted values、range check。\\
一致性 & 跨表、跨字段是否互相矛盾？ & 订单客户编号在客户表中不存在。 & relationships、date order。\\
及时性 & 数据是否足够新？ & 交付数据延迟两周才进入仓库。 & freshness、延迟分布。\\
准确性 & 数据是否反映真实世界？ & 实际交付日期录入错误。 & 抽样核验、外部对账。\\
可追溯性 & 能否追溯来源、规则和版本？ & 指标无法定位输入批次或 SQL 版本。 & lineage、版本记录。\\
\bottomrule
\end{tabular}
\end{table}

质量维度可以交叠。缺失客户编号会影响对象完整性和跨表匹配；超出允许范围的错误日期同时涉及有效性与准确性。按维度描述问题，有助于确定检查方式和修复职责。

\subsection{数据剖析}
\subjectindex{shujupouxi}{数据剖析}

数据剖析通过统计摘要描述结构、缺失、重复和取值分布，为制定质量规则提供依据。剖析报告通常包括以下内容。
\begin{itemize}
  \item 行数与字段数；
  \item 每列缺失数和缺失率；
  \item 每列不同值数量；
  \item 数值字段最小值、最大值和异常值；
  \item 日期字段范围；
  \item 主键重复数和外键不匹配数；
  \item 枚举字段的高频值和未治理新值。
\end{itemize}

剖析结果可用于核对建模假设。例如，若 \texttt{orders.channel} 出现了 Marketplace，而指标字典只允许 Direct、Partner、Online，则需调查 Marketplace 是新渠道、历史错误值，还是上游系统错误。如果它是新渠道，则应更新元数据和指标口径；如果是错误值，则应修复源数据或建立映射规则。

\begin{definition}[数据剖析]
数据剖析是对数据集的结构、字段、缺失、取值分布、键约束、时间范围和跨表关系进行系统摘要的过程，用于了解数据特征并制定检查规则。
\end{definition}

剖析与测试的区别在于：剖析偏探索，测试偏断言。剖析会提出问题，例如“这个字段为什么有四种渠道值”；测试会表达规则，例如“渠道字段只能属于 Direct、Partner、Online”。二者应当配合使用。

\section{缺失与测量误差}

\subsection{缺失机制与均值估计}
\subjectindex{queshijizhi}{缺失机制}
\subjectindex{MCAR}{MCAR}
\subjectindex{MAR}{MAR}
\subjectindex{MNAR}{MNAR}
\subjectindex{xuanzepiancha}{选择偏差}
设 \(Y\) 为真实结果，\(R=1\) 表示该结果进入分析数据。完整记录的均值估计的是 \(\E[Y\mid R=1]\)，而管理者可能想知道 \(\E[Y]\)。只要 \(\Prob(R=1)>0\)，就有恒等式
\[
\E[Y\mid R=1]-\E[Y]
=\frac{\operatorname{Cov}(Y,R)}{\Prob(R=1)}.
\]
将协方差展开为 \(\E[YR]-\E[Y]\E[R]\)，再利用 \(\E[YR]=\Prob(R=1)\E[Y\mid R=1]\)，即可得到该恒等式。偏差由结果与记录机制的协方差以及被观察概率共同决定。

\begin{example}[慢订单更容易缺失]
100 个订单中，80 个普通订单周期均为 2 天，20 个问题订单均为 12 天，真实均值为 4 天。若普通订单全部回写、问题订单只回写 5 个，观察均值为 \((80\times2+5\times12)/85\approx2.59\) 天。完整率达 \(85\%\)，但周期明显被低估。删除缺失记录无法修复这种选择。
\end{example}

完全随机缺失要求缺失与所有相关变量独立；给定已观测协变量后的随机缺失允许缺失依赖这些协变量；非随机缺失则仍依赖未观测结果。实际数据通常不能仅靠一项检验区分后两者。应调查业务记录机制，并用“所有未回写订单均较慢”等情景做敏感性分析。对窗口内已到期但尚未交付的订单，未完成状态本身就是履约证据；已交付订单准时率须与到期承诺口径的 OTD 明确区分（第~\ref{ch:week03-sql-kpi}~章）。

\subsection{经典测量误差与衰减偏差}
\subjectindex{celiangwucha}{测量误差}
设真实线性关系为 \(Y=\alpha+\beta X+\epsilon\)，但只观察到 \(W=X+U\)。若 \(U\) 与 \(X,\epsilon\) 独立、均值为零，且 \(\epsilon\) 与 \(X\) 不相关，使用 \(W\) 回归的总体斜率为
\[
\frac{\operatorname{Cov}(W,Y)}{\operatorname{Var}(W)}
=\beta\frac{\operatorname{Var}(X)}
{\operatorname{Var}(X)+\operatorname{Var}(U)}.
\]
测量噪声使斜率向零衰减。这是经典独立加性误差模型下的结果；若传感器误差与负载相关，方向未必向零。即使数据没有空值、格式全部正确，也可能因为测量偏差而不适合推断。

\subsection{缺失结果的部分识别}
\subjectindex{bufenshibie}{部分识别}
若缺失机制未知，而缺失结果具有已知取值范围，可以求总体均值的上下界。
设共有 \(N\) 个应观察对象，
其中 \(n\) 个被观察，观测值之和为 \(S\)，而每个缺失结果只能确认位于 \([a,b]\)。
则完整总体均值 \(\mu_N\) 必须满足
\begin{equation}
\frac{S+(N-n)a}{N}\le\mu_N\le
\frac{S+(N-n)b}{N}.
\label{eq:missing-bounds}
\end{equation}
推导只需要将每个未知项分别替换为其允许的最小和最大值，再除以总体规模。
该区间是给定已观测值和取值约束所确定的识别区间。它描述相容的总体均值集合，
与抽样不确定性产生的置信区间具有不同含义。

\begin{example}[不完整周期记录的解释]
若应有100张订单，已观察85张，其中80张周期为2天、5张为12天，则已观察周期总和
为220天。假定这批订单已全部结案，遗漏的15张订单经业务核对只知周期在0至20天
之间，则总体平均周期位于 \([2.2,5.2]\) 天。仅报告已观察均值 \(220/85\approx2.59\)
天，会掩盖这个范围。若管理标准是平均周期不超过3天，现有证据不足以判定达标。
\end{example}

上界20天来自对已结案订单的业务核查。若只有非负性条件，均值仍有下界，但通常
没有有限上界。补齐部分遗漏记录或缩小其取值范围，可以收紧识别区间。补充调查
因此可以优先关注对达标判断影响较大的记录。

\subsection{比率误差的精确传播}
\subjectindex{wuchachuanbo}{误差传播}
令真实指标为 \(q=N/D\)，观察值为 \(\tilde q=(N+\delta_N)/(D+\delta_D)\)，其中 \(D>0\) 且 \(D+\delta_D>0\)。相减得
\[
\tilde q-q
=\frac{D\delta_N-N\delta_D}{D(D+\delta_D)}.
\]
若 \(\abs{\delta_D}<D\)，便可给出保守界
\[
\abs{\tilde q-q}
\le\frac{D\abs{\delta_N}+\abs{N}\abs{\delta_D}}
{D(D-\abs{\delta_D})}.
\]
分母小的时候，同样大小的数据错误会造成更大波动。若真实良率为 \(900/1000\)，漏记 100 件不良品使分母变成 900，指标就变成 \(100\%\)。这一变化由漏记不良品造成。

\section{质量规则与统计检验}
\subsection{质量规则的形式化表达}
\subjectindex{zhiliangguize}{质量规则}

设数据表为带有行标识的记录序列 \(R=(r_1,\ldots,r_n)\)。保留行标识才能区分
内容相同的两条记录，避免集合去重掩盖重复问题。质量规则可写为谓词
\[
\phi_j(r_i;R,\mathcal M)\in\{0,1\},
\]
其中 \(\mathcal M\) 包含关联表和规则元数据。单行格式规则只需读取 \(r_i\)，
唯一性需要比较 \(R\) 中其他记录，外键规则则需要关联表。当谓词为0时，
记录 \(r_i\) 违反规则 \(j\)。规则对应的失败行标识集合为
\[
F_j=\{i\in\{1,\ldots,n\}:\phi_j(r_i;R,\mathcal M)=0\}.
\]
若 \(F_j=\varnothing\)，则该规则在当前数据版本上通过；若失败集合不为空，则应记录失败行、规则名称、严重级别、责任人和处理状态。

当 \(n>0\) 时，规则通过比例可定义为
\[
Q_j(R)=1-\frac{|F_j|}{n},
\]
空表应单独报告“无可评估记录”，不将零分母解释为质量满分。该比例仅描述单项规则的符合情况。某些规则只要失败一行就可能导致严重后果，例如主键重复或核心外键缺失；某些规则虽然失败比例不高，但集中在关键客户或高价值订单上，管理影响仍然很大。因此，质量治理应关注失败记录的业务重要性和修复路径。

\begin{convention}[失败记录优先原则]
质量测试返回检验状态和失败记录。记录标识、规则名称、发现时间及异常字段为定位、修复和复核提供依据。
\end{convention}

失败记录表的最小字段可设计为：
\begin{table}[htbp]
\centering
\caption{失败记录表的最小字段}
\begin{tabular}{P{0.22\linewidth}P{0.65\linewidth}}
\toprule
\rowcolor{BDMTableHead}
字段 & 说明\\
\midrule
\texttt{check\_id} & 质量规则的稳定标识，例如 \texttt{unique\_order\_id}。\\
\texttt{dimension} & 对应质量维度，例如 completeness、validity、consistency。\\
\texttt{severity} & error、warning 或 info。\\
\texttt{table} 与 \texttt{column} & 失败发生的数据表和字段。\\
\texttt{row\_id} & 可定位记录的业务键或行标识。\\
\texttt{message} & 面向人阅读的失败解释。\\
\texttt{owner} & 对规则或业务定义负责的角色。\\
\texttt{detected\_at} & 检测时间或数据批次。\\
\bottomrule
\end{tabular}
\end{table}

\subsection{质量规则的程序检验}

数据质量测试可以写成返回违规记录的查询。查询结果同时给出规则是否满足以及需要核查的具体对象。

\begin{table}[htbp]
\centering
\caption{质量规则与失败记录条件}
\begin{tabular}{P{0.20\linewidth}P{0.30\linewidth}P{0.38\linewidth}}
\toprule
\rowcolor{BDMTableHead}
规则类型 & 管理含义 & 失败记录条件\\
\midrule

not null & 必填业务标识存在。 & \texttt{customer\_id is null}\\
unique & 业务对象不被重复计数。 & 同一 \texttt{order\_id} 的记录数大于 1。\\
relationships & 外键能连接到主数据。 & 订单客户在客户表中不存在。\\
accepted values & 字段取值来自约定枚举。 & 渠道不在允许列表中。\\
range check & 数值满足业务边界。 & 订单数量小于或等于 0。\\
date order & 时间顺序符合业务过程。 & 承诺交付日期早于下单日期。\\
freshness & 数据足够新。 & 当前时点距最近一次有效更新的时长超过约定上限。\\
\bottomrule
\end{tabular}
\end{table}

质量规则应同时包含技术表达和业务解释。技术表达说明如何找出失败记录；业务解释说明失败为什么影响指标、模型或行动。例如，\texttt{customer\_id} 缺失不仅是字段为空，它会使订单无法归属区域和客户分层，从而影响区域收入、客户价值和销售责任评价。

\begin{lstlisting}[style=sqlstyle,caption={主键重复失败记录查询}]
SELECT
  order_id,
  COUNT(*) AS row_count
FROM orders
GROUP BY order_id
HAVING COUNT(*) > 1;
\end{lstlisting}

\begin{lstlisting}[style=sqlstyle,caption={客户外键失败记录查询}]
SELECT
  o.order_id,
  o.customer_id
FROM orders AS o
LEFT JOIN customers AS c
  ON o.customer_id = c.customer_id
WHERE o.customer_id IS NOT NULL
  AND c.customer_id IS NULL;
\end{lstlisting}

第一个查询列出重复使用的订单编号，第二个查询列出缺少对应客户主数据的订单。将这些记录与订单金额、客户类别和统计期间关联，即可确定受影响的经营指标及处理优先级。

\section{质量风险与治理资源配置}
\subsection{质量问题对指标的影响}

质量问题对 KPI 的影响可以沿第~\ref{ch:week03-sql-kpi}~章的指标公式追踪。设真实事实集合为 \(O\)，被查询到的事实集合为 \(\tilde O\)。若数据质量问题使得 \(\tilde O\ne O\)，则指标估计
\[
\hat M(\tilde O)
\]
可能偏离真实管理状态 \(M(O)\)。这种偏离有几类常见来源。

\begin{table}[htbp]
\centering
\caption{质量问题对 KPI 的影响路径}
\begin{tabular}{P{0.22\linewidth}P{0.34\linewidth}P{0.30\linewidth}}
\toprule
\rowcolor{BDMTableHead}
质量问题 & 影响机制 & 受影响指标\\
\midrule
订单重复 & 同一事实被重复进入分子或分母。 & 销售额、订单量、延期率。\\
外键缺失 & 内连接丢失事实，或无法分组归属。 & 区域销售额、客户分层 KPI。\\
数量为负或零 & 收入和毛利计算失真。 & 销售额、毛利率。\\
日期顺序错误 & 交付周期和延期判断失真。 & 准时交付率、平均延期天数。\\
枚举不一致 & 分组口径不稳定。 & 渠道收入、渠道准时率。\\
数据延迟 & 当前状态被历史数据替代。 & 库存预警、履约监控。\\
\bottomrule
\end{tabular}
\end{table}

质量问题的严重程度取决于错误类型、影响范围及经营后果。一般字段遗漏可记录后补充；影响财务结算、客户承诺或自动决策的错误，则需根据发布规则暂停使用相关数据并优先修复。

\subsection{质量影响评估}

质量影响评估包括违规规则、具体记录及其对管理判断的影响，需明确失败范围、指标影响、风险等级、处理方式、责任边界及修复证据。

\begin{table}[htbp]
\centering
\caption{质量影响评估模板}
\begin{tabular}{P{0.22\linewidth}P{0.64\linewidth}}
\toprule
\rowcolor{BDMTableHead}
环节 & 需要回答的问题\\
\midrule

失败范围 & 失败记录有多少，集中在哪些客户、产品、区域或时间段？\\
指标影响 & 它会改变哪些指标的分子、分母、分组或时间窗口？\\
风险等级 & 是否影响财务披露、客户承诺、绩效考核或自动化动作？\\
处理方式 & 应阻断发布、带警告发布、补录数据，还是修改规则？\\
责任边界 & 由源系统、主数据团队、指标负责人 还是数据平台团队处理？\\
复盘证据 & 修复后如何证明指标已经恢复可信？\\
\bottomrule
\end{tabular}
\end{table}

例如，\texttt{O107} 的 \texttt{actual\_date} 为空。如果该订单尚未到期且尚未交付，不能仅凭空值认定履约失败；如果已经到期且确认未交付，则应进入到期口径的未准时集合；如果已经交付但系统未回写，则属于数据滞后或集成失败。若来自人工漏填，应补录并检查录入流程。相同的空值在不同业务状态下对应不同的处理，不能由一个统一填充值取代调查。

\subsection{发布决策与预期损失}
\subjectindex{fabukongzhi}{发布控制}
规则失败之后是否停止看板或自动决策，是成本决策。令 \(Z=1\) 表示当前数据存在实质错误，观察到告警证据 \(e\) 后的概率为 \(p=\Prob(Z=1\mid e)\)。若继续发布在错误时损失 \(C_{\rm bad}\)，暂停发布的成本固定为 \(C_{\rm hold}\)，且其他成本可忽略，则
\[
\text{暂停}\quad\Longleftrightarrow\quad
pC_{\rm bad}>C_{\rm hold}.
\]
例如暂停一小时损失 2 万元，错误排程可能损失 100 万元，阈值为 \(p>0.02\)。这些数值为假设参数，实际评价还应区分误报、漏报、恢复时间和已知替代数据。法律或合同要求的强制阻断不应被该简化经济阈值覆盖。

若告警很多，治理投入 \(x_j\) 应按影响而非失败数量分配。设资产 \(j\) 的事故概率近似为 \(p_j(x_j)=p_{j0}\exp(-a_jx_j)\)，损失为 \(L_j\)，预算为 \(B\)，可以求解
\[
\min_{x_j\ge0}\sum_j L_jp_{j0}e^{-a_jx_j}
\quad\text{s.t.}\quad \sum_jx_j\le B.
\]
假定 \(a_j,L_j,p_{j0}>0\)。每增加一单位投入，预期损失的边际下降为
\(a_jL_jp_{j0}e^{-a_jx_j}\)。若两个已获得正投入的对象具有不同边际收益，将少量
预算从收益较低者移向较高者便可降低总损失。因此，内点最优分配要求这些边际收益
相等；未获得投入的对象，其初始边际收益应不高于这一共同值。该条件的约束优化推导见第9章。参数可由历史修复记录或专家评估估计。

\begin{figure}[htbp]
\centering
\begin{tikzpicture}[node distance=0.6cm]
\node[bdmnode] (s) {源事件\\设备与订单};
\node[bdmnode,right=of s] (t) {转换与测试\\规则版本};
\node[bdmnode,right=of t] (k) {指标与模型\\影响范围};
\node[bdmnode,right=of k] (a) {发布与行动\\损失及责任};
\draw[bdmarrow] (s)--(t);
\draw[bdmarrow] (t)--(k);
\draw[bdmarrow] (k)--(a);
\draw[bdmarrow,dashed] (a.north) -- ++(0,0.7) -| (t.north);
\end{tikzpicture}
\caption{血缘追踪与影响评估：上游错误沿依赖关系传播，治理按下游影响排序。}
\end{figure}

\subsection{决策阈值与误差容限}
质量阈值也可以由下游决策反推。假设两条产线的良率估计分别为0.90与0.88，
已知绝对误差界分别为0.02与0.01。它们相容的真值范围为 \([0.88,0.92]\) 与
\([0.87,0.89]\)，区间重叠，因此尚不能保证第一条产线真实良率更高。
这是确定性误差界下的排序结论。

一般地，若 \(\hat k_A-\epsilon_A>\hat k_B+\epsilon_B\)，才能由这些误差界
保证 \(k_A>k_B\)。当两个备选方案的优势很小，质量要求就应更严格；当优势远大于
可能误差，继续追求额外小数位可能没有实际价值。数据质量的误差容限需依据所支持的比较或行动确定。
误差对最优方案的影响可通过敏感性分析与稳健决策研究
（第~\ref{ch:week09-optimization}~章）。

\section{治理组织、契约与事件响应}
\subsection{数据治理与指标治理}

数据治理是组织层面的制度安排，用来确保数据被一致、负责、安全、可追溯地管理。它通常涉及数据负责人、数据管理专员、元数据管理、主数据管理、权限、安全、质量规则、变更流程和审计机制。DAMA-DMBOK 将数据治理放在数据管理知识体系的中心，并围绕其组织数据架构、建模、存储、安全、集成、主数据、元数据、数据质量、数据仓库、BI 和数据科学等领域。

指标治理是数据治理在管理分析中的具体化。KPI 治理包括计算实现、职责分工、变更控制、审计及异常处理：
\begin{enumerate}
  \item 指标负责人：谁对业务定义负责；
  \item 数据管理专员：谁负责数据源、字段含义和质量规则；
  \item 技术实现：SQL、语义层或数据产品中的计算逻辑；
  \item 变更流程：口径变化如何审批、公告和版本化；
  \item 审计记录：每次指标刷新使用了哪些输入数据和代码版本；
  \item 异常处理：质量测试失败时是否阻止发布，或以警告形式发布。
\end{enumerate}

没有治理的指标体系会出现“同名不同义”和“同义不同名”。例如，财务部门的销售额可能以开票为准，运营部门的销售额可能以下单为准，渠道部门的销售额可能以发货为准。三者都有业务理由，但如果不在指标字典中明确，就会在经营会议中制造虚假的分歧。

\begin{convention}[指标治理最小约定]
分析项目中的每个核心 KPI 至少应声明：指标名称、业务定义、事实范围、分子、分母、时间窗口、维度、单位、刷新频率、质量规则、负责人和复核记录。
\end{convention}

\subsection{数据契约、血缘与变更控制}
\subjectindex{shujuqiyue}{数据契约}
\subjectindex{xueyuan}{血缘}
\subjectindex{biangengkongzhi}{变更控制}

数据契约（data contract）是供数方与用数方对数据含义和服务条件的约定，包括字段、粒度、更新频率、允许取值、变更方式及通知程序。约定可转化为自动检查规则，并据变更记录核对上下游一致性。

\begin{table}[htbp]
\centering
\caption{数据契约的最小内容}
\begin{tabular}{P{0.22\linewidth}P{0.63\linewidth}}
\toprule
\rowcolor{BDMTableHead}
内容 & 说明\\
\midrule

数据对象 & 表、视图或事件流的名称与业务含义。\\
粒度 & 一行数据代表什么业务事实。\\
字段定义 & 字段类型、单位、空值含义和允许取值。\\
更新约定 & 刷新频率、最大延迟和补数方式。\\
质量规则 & 必填、唯一、外键、枚举、日期顺序等约束。\\
变更流程 & 字段删除、类型变化、枚举新增如何通知与审批。\\
责任人 & 上游负责人、数据管理专员和下游使用方。\\
\bottomrule
\end{tabular}
\end{table}

血缘记录回答“指标从哪里来”的问题。一个可信指标至少应能追溯到源表、转换脚本、质量测试、数据批次和发布时间。变更控制回答“口径为什么变”的问题。没有血缘和变更控制，即使当前数值正确，组织也难以解释历史对比中的差异。

\subsection{发布控制与问题分级}

质量检查需要行动规则。若所有失败都被视为同等严重，团队会陷入告警疲劳；若所有失败都只是记录日志，关键指标可能在错误数据上发布。按发布策略，质量规则可分为三个级别：

\begin{table}[htbp]
\centering
\caption{质量规则严重级别}
\begin{tabular}{P{0.16\linewidth}P{0.38\linewidth}P{0.34\linewidth}}
\toprule
\rowcolor{BDMTableHead}
级别 & 发布策略 & 示例\\
\midrule

Error & 阻止关键表或关键 KPI 发布。 & 主键重复、核心外键缺失、关键日期非法。\\
Warning & 允许发布但必须标记并跟踪。 & 新渠道值、少量非核心字段缺失。\\
Info & 记录趋势，用于后续改进。 & 字段分布变化、缺失率轻微波动。\\
\bottomrule
\end{tabular}
\end{table}

问题级别依据业务影响确定。例如，普通订单的地址字段缺失可能只是 warning；战略客户订单的实际交付日期缺失可能应升级为 error，因为它直接影响 SLA、客户沟通和绩效考核。

发布控制可以表达为函数
\[
G(D,\Phi)=
\begin{cases}
\text{block}, & \exists j: F_j\ne \varnothing \text{ 且 } s_j=\text{error},\\
\text{warn}, & \text{无 error 失败，且存在 warning 失败},\\
\text{pass}, & \text{otherwise}.
\end{cases}
\]
其中 \(s_j\) 是规则 \(j\) 的严重级别。实际系统还会考虑失败比例、失败金额、客户等级和是否处于财务关账期等条件。

\subsection{质量事件响应}
\subjectindex{shijianxiangying}{事件响应}

当发布控制失败时，组织需要一套可重复的事件响应流程。最低流程包括：
\begin{enumerate}
  \item 识别：质量测试产生失败记录，并记录数据批次；
  \item 定级：根据指标影响、客户影响和自动化风险确定严重级别；
  \item 隔离：暂停受影响指标发布，或在看板上显式标注警告；
  \item 归因：定位问题发生在源系统、主数据、管道还是指标逻辑；
  \item 修复：补录、回填、重跑、修改规则或更新数据契约；
  \item 复盘：记录根因、影响范围、修复证据和预防措施。
\end{enumerate}

质量事件响应持续记录问题发现、处置与修复情况。分析项目中的质量报告可以简化，但必须保留失败记录、影响解释和处理建议。

\section{数据目录与质量监测的企业应用}

\subsection{Collibra 与 Froedtert 的数据目录}
\subjectindex{Collibra}{Collibra}
\subjectindex{Froedtert}{Froedtert}
Collibra 的 Froedtert 与 Medical College of Wisconsin 案例介绍了目录建设与数据治理实践，组织内部将相关平台称为 Data 360\cite{collibra04,collibraGovernance}。案例的分析对象是指标定义、数据发现与责任安排。

该机构已有电子病历、财务与人力资源记录，数据治理面对系统分散、指标定义差异和责任安排等问题。其组织形式结合中心协调与业务参与，通过指标审批确认定义和负责人。Data 360 保存数据说明等元数据，敏感明细的访问另由相应权限管理。

例如，预约等待天数可以从提出预约请求起算，以最早可预约日期或患者实际选择日期为终点。两种定义分别描述服务供给和患者实际等待。明确用途后，指定定义负责人并记录修改程序，可以使不同部门使用相同的计算依据。

在封测算例中，“良品”需分别按首检、返工后终检和客户验收三个环节定义。业务定义负责人确认各项口径，数仓实现采用对应定义；变更单记录生效时间、历史重算规则及下游通知情况。

这一安排使首检良率、最终良率与客户验收通过率分别服务于过程控制、成品评价和交付评价。定义确认之后，还需依据第~\ref{ch:week02-enterprise-data-model}~章的激励分析，审查各指标进入考核时的行为后果。

\subsection{Monte Carlo 与 Assurance 的数据质量监测}
\subjectindex{MonteCarlo}{Monte Carlo（企业）}
\subjectindex{Assurance}{Assurance}
\subjectindex{keguancexing}{可观测性}
Monte Carlo 的 Assurance 案例披露了在业务增长中引入数据可观测性、改善数据可信度的实践\cite{monte04}。其 2025 年非结构化数据监控公告还将范围扩展到文档等资产\cite{unstructured04}。字段规则检验与运行过程监测承担不同职责。

字段测试判断预先声明的条件是否成立，数据可观测性则持续监测数据量、分布、更新时间、模式及血缘变化。准时到达的文件可能重复昨日内容，此时需要检查记录时间和内容标识。市场需求骤降也会改变数据分布，应结合来源系统与经营情况区分真实变化和处理错误。

质量事件可用平均检测时间 MTTD 和平均修复时间 MTTR 描述，同时记录错误发布的持续时间。设每小时错误决策损失为 \(\ell\)，事件持续时间为 \(T\)，则相应损失为 \(\ell T\)。检测改进缩短发现之前的时间，责任分配和修复资源则影响发现之后的处置时间，两者共同决定总损失。

\subsection{数据目录、治理流程与敏感信息识别}
\subjectindex{Alation}{Alation}
\subjectindex{BigID}{BigID}
Alation 数据目录的公开产品资料介绍了数据资产发现、业务术语、来源说明和血缘
信息的组织方式\cite{alationCatalog2026}。对分析人员而言，目录的作用是缩短
查找和理解数据的时间；对治理人员而言，它提供责任人、规则和质量状态的维护
位置。目录收录了一张订单表，并不自动说明订单金额是否含税或其每日更新时间，
这些含义仍须由业务与技术负责人填写并核对。

Collibra 平台的产品范围包括目录、血缘、质量监测、隐私及人工智能治理
\cite{collibraPlatform2026}。其治理流程可以把“准时交付率”的定义变更分配给
指定负责人，记录审批与生效日期，再标识受影响的报告。BigID 着重于跨结构化
与非结构化数据发现、敏感信息分类及风险处置\cite{bigidDiscovery2026}；
当工艺文档或客户合同进入检索系统时，应先识别其中的个人信息、商业秘密和授权
范围。这三类机制分别解决发现与理解、定义与责任、敏感内容识别等问题；
实际平台能力和可用接口应以部署环境中的测试结果确认。

以“让智能体解释客户 C17 的延期”为例，目录应能找到订单与批次表及其指标定义，
血缘记录说明每日看板使用的源表版本，质量检查标记缺失的承诺日期，敏感信息
分类表明客户合同可供哪些角色检索。查询服务应根据当前用户身份限制返回行，
文档检索应执行同一授权边界。只有将记录规则、权限规则和使用记录分别核查，
目录中的“可信”标记才具有明确的业务含义。

\section{检索增强生成中的数据质量}
\subjectindex{zhengjuyouxiaoqi}{证据有效期}
RAG 将检索到的材料作为模型回答的上下文。因此质量检查除了表字段，还应包含文档有效期、版本、来源、权限、切片边界和引用定位。旧版工艺规则即使文本提取完全正确，也可能不适合今天的排程。扫描表格中丢失单位，会把正确的数字转成错误的约束。

对每个证据片段 \(d\)，可保存
\[
(\text{id},\text{source},\text{version},
\text{valid\_from},\text{owner},\text{scope},\text{hash}).
\]
哈希标识用于检测内容变化，来源真实性由发布主体和签发记录核查。访问权限在检索阶段执行，撤销或过期材料同步更新其索引与缓存状态，使模型获得授权范围内的有效文档。

检索召回率与答案正确率需要分别测量。正确答案不存在于知识库时，应允许“证据不足”；若答案依赖两份相互矛盾的规范，应显示冲突和适用时间。回答应列出对应的有效来源及引用位置。

\section{数据质量检查实例}

配套样例提供 \path{data/sample/orders_quality_issues.csv}，其中包含若干刻意构造的问题：重复订单号、缺失客户编号、未知客户、未知产品、无效渠道、无效优先级、非正数量、承诺日期早于下单日期、实际交付日期缺失。实验任务是使用 Python 或 SQL 找出违规记录，并分析每类违规对 KPI 的影响。

配套代码由四个模块组成：
\begin{itemize}
  \item \path{src/bdm_decision/quality/data_profile.py}：生成字段级剖析信息；
  \item \path{src/bdm_decision/quality/checks.py}：提供 not null、unique、accepted values、relationships、positive number 和 date order 等检查；
  \item \path{src/bdm_decision/quality/validation_rules.py}：组合订单样例数据所需的规则；
  \item \path{src/bdm_decision/cases/week04_data_quality.py}：加载样例数据并生成质量报告。
\end{itemize}

最小运行方式为：
\begin{lstlisting}[style=pythonstyle,caption={运行第4章数据质量案例}]
import bdm_decision.cases.week04_data_quality as week04

print(week04.render_quality_report_markdown(dirty=True))
\end{lstlisting}

实验输出包括字段摘要、违规记录、严重级别统计和质量维度统计。每类违规记录还应说明受影响的指标、可能的数值偏差及处理方法。

\begin{table}[htbp]
\centering
\caption{样例数据中的典型失败}
\begin{tabular}{P{0.20\linewidth}P{0.28\linewidth}P{0.36\linewidth}}
\toprule
\rowcolor{BDMTableHead}
失败类型 & 样例 & 管理解释\\
\midrule

重复订单 & \texttt{O103} 重复。 & 销售额、订单数和准时率可能被重复计入。\\
未知客户 & \texttt{C999} 不在客户表中。 & 订单无法归属区域和客户分层。\\
未知产品 & \texttt{P999} 不在产品表中。 & 价格、成本和毛利无法可靠计算。\\
非正数量 & \texttt{O102} 数量为 0，\texttt{O109} 数量为负。 & 收入和毛利计算失真。\\
日期顺序错误 & \texttt{O105} 承诺日期早于下单日期。 & 交付周期和 SLA 判断失真。\\
未治理枚举 & Marketplace、Urgent。 & 渠道和优先级分组口径不稳定。\\
\bottomrule
\end{tabular}
\end{table}

\subsection{质量报告的核查}

一个可审计质量报告应当满足五项最低要求：
\begin{enumerate}
  \item 说明数据集、批次、行数和字段数；
  \item 列出质量规则、维度和严重级别；
  \item 保留违规记录及其数量；
  \item 说明失败对指标和管理行动的影响；
  \item 记录处理状态、责任人和复核入口。
\end{enumerate}

在企业实践中，质量报告还应与任务调度、指标发布、告警系统和工单系统连接。实验报告可采用 Markdown、CSV 或 notebook，须保留可复现的违规记录及其管理影响。

\section*{本章小结}
\addcontentsline{toc}{section}{本章小结}

数据质量是相对于具体用途的要求。数据剖析描述记录特征，规则检查发现约束违反，对账和抽样复核评价记录准确性。缺失机制、测量误差和分母误差分别通过不同途径影响统计量，需要采用相应的偏差分析、识别区间或误差界。

质量管理据此比较错误发布的预期损失、暂停成本与修复收益，并分配治理资源。数据目录、质量工作和敏感信息发现各有职责，要通过实际数据来源与授权检查它们的覆盖范围。文档型知识库还需要检查来源、有效期与访问范围。统计摘要与图形对经营特征的描述，同样以数据来源及质量条件明确为前提。

\chapterreading{Wang 与 Strong 从数据使用者视角讨论超出准确性的质量维度，Pipino 等讨论
质量评价的组织与测量问题\cite{wang1996,pipino2002}。Redman、English 与 Olson
提供数据质量管理的系统论述\cite{redman1996,english1999,olson2003}。
ISO/IEC 25012 与 DAMA-DMBOK 分别提供质量模型和治理知识结构
\cite{iso25012,dama2017}。dbt 的数据测试文档可作为将声明式规则转化为可执行
检查的技术入口\cite{dbt-tests}。数据目录、治理流程与敏感数据发现的
产品机制可分别查阅\cite{alationCatalog2026,collibraPlatform2026,bigidDiscovery2026}。工具规则的通过范围仍须与业务真实性区分。}

\section*{思考与习题}
\addcontentsline{toc}{section}{思考与习题}

\begin{exercise}[从完整率到偏差]
在慢订单算例中，令问题订单回写比例为 \(r\in[0,1]\)，写出观察均值随 \(r\) 的函数。解释为什么仅要求完整率超过 \(80\%\) 无法保证均值误差足够小。
\end{exercise}
\begin{exercise}[发布与检索控制]
某指标错误概率为 \(0.04\)，错误行动损失为 80，暂停成本为 2，单位均为万元。求最优发布选择，并解释损失估计不确定时如何做敏感性分析。再为一份过期设备 SOP 设计检索前的有效性检查。
\end{exercise}

\begin{exercise}
\begin{enumerate}
  \item 对 \path{data/sample/orders_quality_issues.csv} 做字段剖析，报告每列缺失数、缺失率、不同值数量和日期范围。
  \item 写出检查 \texttt{order\_id} 唯一性的 SQL，并说明重复订单会如何影响销售额和准时交付率。
  \item 写出检查订单客户外键的 SQL，并说明外键缺失会如何影响区域级 KPI。
  \item 将“渠道取值必须属于 Direct、Partner、Online”写成质量规则，并为 Marketplace 取值设计处理方案。
  \item 为 \(\mathrm{on\_time\_delivery\_rate}\) 设计一组发布控制，区分 error、warning 和 info。
  \item 选择一个失败记录，写出其质量维度、影响指标、责任角色、处理方式和复盘证据。
\end{enumerate}
\end{exercise}

\begin{exercise}[缺失范围与补充调查]
在式~\eqref{eq:missing-bounds}~的例子中，若复核发现遗漏的15张订单中有5张周期均为
12天，其余仍只知处于0至20天之间，重新计算总体均值范围。若业务要求平均周期
不超过3天，这次复核是否足以改变判断？解释为什么这个范围不应被标为95\%置信区间。
\end{exercise}

\begin{exercise}[把质量预算投向决策敏感位置]
两条产线的估计良率为0.90和0.88，误差界分别为0.02和0.01。
假定只能缩小其中一个误差界，求保证第一条产线良率更高所需的条件，并说明在哪一种
选择下严格排序仍不可能得到保证。再讨论误差界若只是未经验证的经验判断，结论应
如何表述。不要把误差界、标准误和质量规则失败率视为同一个量。
\end{exercise}
